\itemhead{2.2 思考 · 未编号题组（教材p.45）}{XJ-B1-C02-TH-P045-Y0479}
\label{XJ-B1-C02-TH-P045-Y0479}
{\footnotesize 湘教版教材p.45 / PDF第52页 / 2.2 思考 / 未编号题组}
\field{原题}{在初中，我们探讨过一次函数与二元一次方程的内在联系。对于二次函数与一元二次方程，它们之间又有怎样的联系呢？}
\field{参考答案}{方程 $ax^2+bx+c=0$ 的实根是二次函数图象与 $x$ 轴交点的横坐标；交点个数由判别式 $\Delta=b^2-4ac$ 的符号决定（$a\ne0$）。}
{\small \field{关联知识}{Q01 二次函数、零点、判别式与根系关系（教材帮p.65 / PDF74）}
\field{方法与易错点}{按已知零点、顶点或系数选择交点式、顶点式或一般式；根平方和用根系关系消元。 易错：遗漏 $a
e0$；把二重根算成两个不同零点；有两正根只查积正。}
\field{答案依据}{答案册未提供本栏目的独立对应答案；依据原题、教材概念补写参考回应。开放性总结不设唯一数值答案。}}
